% This contents of this file will be inserted into the _Solutions version of the
% output tex document.  Here's an example:

% If assignment with subquestion (1.a) requires a written response, you will
% find the following flag within this document: <SCPD_SUBMISSION_TAG>_1a
% In this example, you would insert the LaTeX for your solution to (1.a) between
% the <SCPD_SUBMISSION_TAG>_1a flags.  If you also constrain your answer between the
% START_CODE_HERE and END_CODE_HERE flags, your LaTeX will be styled as a
% solution within the final document.

% Please do not use the '<SCPD_SUBMISSION_TAG>' character anywhere within your code.  As expected,
% that will confuse the regular expressions we use to identify your solution.

\def\assignmentnum{1 }
\def\assignmenttitle{XCS234 Assignment \assignmentnum}
\input{macros}
\begin{document}
\pagestyle{myheadings} \markboth{}{\assignmenttitle}

% <SCPD_SUBMISSION_TAG>_entire_submission

This handout includes space for every question that requires a written response.
Please feel free to use it to handwrite your solutions (legibly, please).  If
you choose to typeset your solutions, the |README.md| for this assignment includes
instructions to regenerate this handout with your typeset \LaTeX{} solutions.
\ruleskip

\LARGE
1.f
\normalsize

% <SCPD_SUBMISSION_TAG>_1f
\begin{answer}
  % ### START CODE HERE ###
  % \usepackage{color}
  % \usepackage{tabularray}
  \definecolor{Shark}{rgb}{0.121,0.137,0.156}
  \begin{table}
  \centering
  \caption{Discount factor and corresponding expected returns for transitions to "left" (a=0) action selection at initial state (s=0)}
  \begin{tblr}{
    cell{1}{2} = {c=2}{fg=Shark},
    cell{1}{4} = {c=2}{},
  }
  Current strength & transition~$/gamma$ &             & $(q_{*}(s=0, a=0), q_{*}(s=0, a=1))$ &                  \\
                  & policy iter.        & value iter. & policy iter.                           & value iter.      \\
  WEAK             & 0.65                & 0.65        & (0.0143, 0.0137)                       & (0.0143, 0.0137) \\
  MEDIUM           & 0.76                & 0.76        & (0.0208, 0.0203)                       & (0.0208, 0.0203) \\
  STRONG           & 0.92                & 0.92        & (0.0599, 0.0582)                       & (0.0599, 0.0582) 
  \end{tblr}
  \end{table}

  The table above contains the values of \gamma (to 2 decimal places) for which the optimal agent transitioned from selecting "right" at 
  the initial state, to when it instead selected "left" instead. These \gamma values of transition were calculated for both policy 
  iteration and value iteration dynamic progamming algorithms across current strength settings WEAK, MEDIUM and STRONG of the Markov 
  Process RiverSwim.

  From the table, both policy and value iteration methods have identical transition discount factors for the three different 
  current strengths, as well as identical action-value function returns that the optimal agent compared to favor left over right at the far-left 
  initial state \s_{1}. Stronger current decreases the probability of successful swim right attempts in the
  stochastic dynamical state transition matrix \T. Transition discount factors monotonically increased from WEAK to STRONG currents 
  as the optimal agent required more foresight into expected future return (incentive) to overcome the reduced probability of 
  success of actually swimming right after deciding to swim right.
\clearpage


\LARGE
2.a
\normalsize

% <SCPD_SUBMISSION_TAG>_2a
\begin{answer}
  % ### START CODE HERE ###
  \definecolor{Shark}{rgb}{0.121,0.137,0.156}
  \begin{table}
  \centering
  \caption{Discount factor and corresponding expected returns for transitions to "left" (a=0) action selection at initial state (s=0)}
  \begin{tblr}{
    cell{1}{2} = {c=2}{fg=Shark},
    cell{1}{4} = {c=2}{},
  }
  Horizon \H  & Expected Return & Execution sequence under optimal policy, \gamma = 1      \\
              &                 & (action, state transition, resultant reward) \\
  4           & 3               & (sell, 3 \leftarrow 2, +1), (sell, 2 \leftarrow 1, +1), (sell, 0 \leftarrow 1, +1), (buy, 0 \leftarrow 1, 0) \\
  4           & 3               & (buy, 3 \leftarrow 4, 0), (sell, 4 \leftarrow 3, +1), (sell, 3 \leftarrow 2, +1), (sell, 2 \leftarrow 1, +1) \\
  5           & 4               & (sell, 3 \leftarrow 2, +1), (sell, 2 \leftarrow 1, +1), (sell, 0 \leftarrow 1, +1), (buy, 0 \leftarrow 1, 0), (sell, 0 \leftarrow 1, +1) \\
  6           & 4               & (sell, 3 \leftarrow 2, +1), (sell, 2 \leftarrow 1, +1), (sell, 0 \leftarrow 1, +1), (buy, 0 \leftarrow 1, 0), (sell, 0 \leftarrow 1, +1), (buy, 0 \leftarrow 1, 0) \\
  7           & 100             & (buy, 3 \leftarrow 4, 0), (buy, 4 \leftarrow 5, 0), (buy, 5 \leftarrow 6, 0), (buy, 6 \leftarrow 7, 0), (buy, 7 \leftarrow 8, 0), (buy, 8 \leftarrow 9, 0), (buy, 9 \leftarrow 10, 100) \\
  8           & 100             & (buy, 3 \leftarrow 4, 0), (buy, 4 \leftarrow 5, 0), (buy, 5 \leftarrow 6, 0), (buy, 6 \leftarrow 7, 0), (buy, 7 \leftarrow 8, 0), (buy, 8 \leftarrow 9, 0), (buy, 9 \leftarrow 10, 100), (sell, 10 \leftarrow 10, 0) \\
  9           & 101             & (sell, 3 \leftarrow 2, +1), (buy, 2 \leftarrow 3, 0), (buy, 3 \leftarrow 4, 0), (buy, 4 \leftarrow 5, 0), (buy, 5 \leftarrow 6, 0), (buy, 6 \leftarrow 7, 0), (buy, 7 \leftarrow 8, 0), (buy, 8 \leftarrow 9, 0), (buy, 9 \leftarrow 10, 100) \\
  10          & 101             & < same as above for sequence entries 1 through 8 >, (buy, 9 \leftarrow 10, 100), (sell, 10 leftarrow 10, 0) \\
  11          & 102             & (sell, 3 \leftarrow 2, +1),  (sell, 2 \leftarrow 1, +1), (buy, 1 \leftarrow 2, 0), (buy, 2 \leftarrow 3, 0), (buy, 3 \leftarrow 4, 0), (buy, 4 \leftarrow 5, 0), (buy, 5 \leftarrow 6, 0), (buy, 6 \leftarrow 7, 0), (buy, 7 \leftarrow 8, 0), (buy, 8 \leftarrow 9, 0), (buy, 9 \leftarrow 10, 100) \\
  12          & 102             & < same as above for sequence entries 1 through 10 >, (buy, 9 \leftarrow 10, 100), (sell, 10 leftarrow 10, 0) \\
  \end{tblr}
  \end{table}
  An optimal policy is one that maximizes expected cumulative future return from the current time step to the horizon.  While the expected
  cumulative future return is unique in Markov Decision Processes, there may be multiple distinct policies that generate it; that is, an 
  optimal policy from policy or value iteration is not necessarily unique.
  
  For the Inventory MDP in this problem, an optimal policy can change dramatically depending on the number of steps that remain in the episode during 
  execution before the horizon is reached.  Notice in the above table that expected return jumps from 4, for \H = 6, to 100 for \H = 7.  Also,
  for the same two horizon settings execution sequences transistioned from having both buy and sell actions when \H = 6, to buy only action 
  when \H = 7.

Inspection of the table reveals that execution sequences under the optimal policy contain both buy and sell actions when 4 \lte H \lte 6 and when H \gte 8.
  % ### END CODE HERE ###
\end{answer}
% <SCPD_SUBMISSION_TAG>_2a
\clearpage

\LARGE
2.b
\normalsize

% <SCPD_SUBMISSION_TAG>_2b
\begin{answer}
  % ### START CODE HERE ###
  Settin the horizon \H \gte 7 will result in an execution sequence under the optimal policy
  that terminates with a fully stocked inventory.  See table. Notice that a fully stocked inventory
  occurs regardless of whether the agent performs action selection in the terminal state s_{10} or
  only enters into it without performing actions in that state.
  % ### END CODE HERE ###
\end{answer}
% <SCPD_SUBMISSION_TAG>_2b
\clearpage

\LARGE
2.c
\normalsize

% <SCPD_SUBMISSION_TAG>_2c
\begin{answer}
  % ### START CODE HERE ###
  In an infinite-horizon scenario with discount factor setting \gamma = 0:
  1. When starting in state s = 3, the sequence of actions that the agent takes under the optimal
  policy is (sell, 3 \leftarrow 2, +1), (sell, 2 \leftarrow 1, +1), (sell, 1 \leftarrow 0, +1), (sell, 0 \leftarrow 0, 0), (sell, 0 \leftarrow 0, 0), ...
  ad infinitum.  The sequence cannot terminate as it never exits state transition 0 \leftarrow 0 loop 
  caused by perpetual greedy action selection--the consequence of a myopic agent via setting \gamma = 0.

  2. Similarly, for the scenario starting as s = 9, the execution squence under the optimal policy migrates
  to the edge of the state space, in this case s = 10--the terminal state--due to the inability of the agent 
  to consider non-local return possibilities.  The execution sequence is: (buy, 9 \leftarrow 10, 100), (sell, 10 \leftarrow 10, 0)
  % ### END CODE HERE ###
\end{answer}
% <SCPD_SUBMISSION_TAG>_2c
\clearpage

\LARGE
2.d
\normalsize

% <SCPD_SUBMISSION_TAG>_2d
\begin{answer}
  % ### START CODE HERE ###
  In the infinite-horizon scenario with reward discounting, it is possible 
  to choose a fixed value \gamma \in \left[0, 1\right) such that the optimal 
  policy never fully stocks the inventory. Setting \gamma to a value for which 
  the cummulative expected discounted return for selling does not surpass that for buying at
  state s = 3 would result in the inequality 100 \cdot \gamma^{6} \lt 1 + 1 \cdot \gamma + 1 \cdot \gamma^2.
  Numerically estimating the real positive root for this polynomial equation we find \gamma \lt 0.5106.
  Testing about this transition \gamma value under value iteration with tolerance 1e-3, we find that
  the action selected at state s = 3 with discount factor set to \gamma=0.52 is sell--which would end with the inventory
  fully stocked. Setting \gamma=0.51 results in selected action sell whereby the agent's policy would 
  never fully stock the inventory at the end of the day.
  % ### END CODE HERE ###
\end{answer}
% <SCPD_SUBMISSION_TAG>_2d
\clearpage


\LARGE
3.a
\normalsize

% <SCPD_SUBMISSION_TAG>_3a
\begin{answer}
  % ### START CODE HERE ###
  % ### END CODE HERE ###
\end{answer}
% <SCPD_SUBMISSION_TAG>_3a
\clearpage


\LARGE
3.b
\normalsize

% <SCPD_SUBMISSION_TAG>_3b
\begin{answer}
  % ### START CODE HERE ###
  % ### END CODE HERE ###
\end{answer}
% <SCPD_SUBMISSION_TAG>_3b
\clearpage

% <SCPD_SUBMISSION_TAG>_entire_submission

\end{document}
